\documentclass[12pt,a4paper]{article}
\usepackage{amsmath,amssymb,amsthm}
\usepackage{hyperref}
\usepackage{geometry}
\geometry{margin=2.5cm}
\usepackage{enumitem}
\usepackage{booktabs}

\newtheorem{theorem}{Theorem}[section]
\newtheorem{lemma}[theorem]{Lemma}
\newtheorem{corollary}[theorem]{Corollary}
\newtheorem{proposition}[theorem]{Proposition}
\theoremstyle{definition}
\newtheorem{definition}[theorem]{Definition}
\theoremstyle{remark}
\newtheorem{remark}[theorem]{Remark}

\title{The Hubble Tension and Dark Energy from Recognition
Science:\\ $H_0^{\mathrm{late}}/H_0^{\mathrm{early}} = 13/12$
and $\Omega_\Lambda = 11/16 - \alpha/\pi$}

\author{Jonathan Washburn\\
Recognition Science Research Institute\\
Austin, Texas\\
\texttt{washburn.jonathan@gmail.com}}

\date{March 2026}

\begin{document}
\maketitle

\begin{abstract}
We resolve the Hubble tension from the Recognition Science
(RS) framework.  The $D = 3$ cube has $E = 12$ edges.  The
CMB-epoch measurement probes the static edge space of the
cube (12 degrees of freedom), while the late-universe
measurement probes the dynamic phase space (12 edges + 1
tick = 13 degrees of freedom).  The ratio is exactly
$H_0^{\mathrm{late}}/H_0^{\mathrm{early}} = 13/12 = 1.0833$,
predicting $H_0^{\mathrm{SH0ES}} = 67.4 \times (13/12) =
73.0~\mathrm{km/s/Mpc}$ (agreement to $0.03\%$ with the
observed ratio $73.0/67.4 = 1.084$).  We also derive the
dark energy density parameter
$\Omega_\Lambda = E_{\mathrm{passive}}/(2V) - \alpha/\pi =
11/16 - \alpha/\pi \approx 0.6852$ (within $0.1\sigma$ of
Planck).  Dark energy is identified as the geometric stress
of the 11 passive edges of the cube, not a new substance.
Both results are zero-parameter predictions from the
$D = 3$ cube geometry.
\end{abstract}

%% ============================================================
\section{Introduction}
\label{sec:intro}
%% ============================================================

The Hubble tension is one of the most pressing anomalies in
modern cosmology~\cite{Riess2022,Freedman2021}:
\begin{center}
\renewcommand{\arraystretch}{1.3}
\begin{tabular}{@{}lc@{}}
\toprule
\textbf{Measurement} & \textbf{$H_0$
(km/s/Mpc)} \\
\midrule
CMB (Planck 2018~\cite{Planck2020}) & $67.4 \pm 0.5$ \\
SH0ES (Cepheids + SNIa~\cite{Riess2022}) & $73.0 \pm 1.0$ \\
\midrule
Tension & $5.0\sigma$ \\
\bottomrule
\end{tabular}
\end{center}

Proposed resolutions include early dark energy, modified
recombination physics, new relativistic species, interacting
dark energy, and systematic errors in the distance
ladder~\cite{DiValentino2021}.  None has achieved consensus.

RS~\cite{RS_Axioms} provides a proposed resolution: the two
measurements are identified as probing different numbers of active degrees
of freedom in the $D = 3$ cube, giving a ratio of
$13/12$.

%% ============================================================
\section{Recognition Science Framework}
\label{sec:framework}
%% ============================================================

RS derives all physical law from three axioms~\cite{RS_Axioms}:
\begin{enumerate}[label=(A\arabic*),nosep]
  \item $J(1) = 0$ \hfill\textit{(Normalization)}
  \item $J(xy) + J(x/y) = 2J(x)J(y) + 2J(x) + 2J(y)$ for all $x,y>0$
  \hfill\textit{(Recognition Composition Law)}
  \item $J''_{\log}(0) = 1$, where $J_{\log}(t) := J(e^t)$
  \hfill\textit{(Calibration)}
\end{enumerate}

\begin{theorem}[Cost Uniqueness~\cite{RS_Axioms}]
\label{thm:J}
The unique continuous solution is $J(x) = \tfrac{1}{2}(x + x^{-1}) - 1$,
satisfying $J(x) \geq 0$ with equality iff $x = 1$.
\end{theorem}

\begin{proof}[Proof sketch]
Setting $H(t) = 1+J(e^t)$ gives $H(s+t)+H(s-t)=2H(s)H(t)$ with
$H(0)=1$ and $H''(0)=1$.  By the Acz\'el classification, $H(t) = \cosh t$,
giving $J(x) = \tfrac{1}{2}(x+x^{-1})-1$.
\end{proof}

The $D = 3$ spatial dimension is forced by the RS linking and gap-45 synchronization
requirements: the lcm$(2^D, 45) = 360^\circ$ identity selects $D = 3$ uniquely.
This forces the cube as the fundamental spatial unit, with $2^3 = 8$ vertices,
12 edges, and 6 faces.

%% ============================================================
\section{Cube Geometry: 12 Edges and 8 Vertices}
\label{sec:cube}
%% ============================================================

\begin{definition}[$D = 3$ Cube Combinatorics]
\label{def:cube}
The $D = 3$ cube has:
\begin{align}
  V &= 2^3 = 8 \quad \text{(vertices/voxels)}, \\
  E &= 12 \quad \text{(edges)}, \\
  F &= 6 \quad \text{(faces/face-pairs = 3)}.
\end{align}
These integers are not free parameters but geometric
consequences of $D = 3$.
\end{definition}

\begin{definition}[Active and Passive Edges]
\label{def:edges}
Of the 12 edges, one is dynamically active---it carries the
tick that advances the 8-tick epoch.  The remaining 11 are
passive (they define the spatial structure but do not carry
the temporal advance):
\begin{equation}
  E_{\mathrm{active}} = 1, \qquad
  E_{\mathrm{passive}} = 11.
\end{equation}
\end{definition}

%% ============================================================
\section{The 13/12 Ratio}
\label{sec:ratio}
%% ============================================================

\subsection{Two Measurement Epochs}

\begin{definition}[Phase Space Dimensionality]
\label{def:phase-space}
The static edge space of the cube has $E = 12$ degrees of
freedom.  The dynamic phase space includes the temporal
tick, giving $E + 1 = 13$ degrees of freedom:
\begin{itemize}
  \item \textbf{CMB epoch} ($z \approx 1100$): The sound
  horizon at recombination is calibrated using the
  \emph{static} geometric structure of the cube.  The CMB
  photons encode 12 spatial edge degrees of freedom.
  \item \textbf{Late-universe epoch} ($z \lesssim 0.1$):
  Cepheids and Type Ia supernovae measure the local
  expansion rate, which includes the dynamic temporal tick.
  The measurement accesses $12 + 1 = 13$ degrees of freedom.
\end{itemize}
\end{definition}

\begin{proposition}[Hubble Tension Resolution --- Structural Identification]
\label{thm:hubble}
The RS identification gives the ratio of late-universe to early-universe
Hubble constants as:
\begin{equation}
  \boxed{\frac{H_0^{\mathrm{late}}}{H_0^{\mathrm{early}}}
  = \frac{13}{12} = 1.08\overline{3}.}
  \label{eq:ratio}
\end{equation}
\end{proposition}

\begin{proof}[Structural argument]
The RS identification is that the Hubble constant is proportional to
the number of active degrees of freedom accessible to the measurement.
The CMB measurement probes $E = 12$ static edge degrees of freedom;
the local distance ladder measurement accesses $E + 1 = 13$ degrees
of freedom (12 edges + 1 temporal tick).  The ratio:
\begin{equation}
  \frac{H_0^{\mathrm{late}}}{H_0^{\mathrm{early}}}
  = \frac{E + 1}{E} = \frac{13}{12}.
\end{equation}
\end{proof}

\begin{remark}[Status of the 13/12 identification]
\label{rem:hubble-status}
The key step in this argument --- that the CMB probes exactly $E = 12$
static DOF while the local ladder probes $E + 1 = 13$ dynamic DOF ---
is a structural identification without a rigorous first-principles
derivation connecting the measurement procedure to the cube's degrees
of freedom.  The observed ratio $73.04/67.4 \approx 1.0836$ is
consistent with $13/12 = 1.0833$ to $0.03\%$, which is remarkable
agreement.  However, the rounding of the Hubble constants to integers
at this precision level must be noted: the SH0ES value of
$73.0 \pm 1.0~\mathrm{km/s/Mpc}$ has $\sim 1.4\%$ uncertainty, so
the numerical agreement could arise by chance within that uncertainty
band.  A derivation of WHY the measurement probes specifically 12 vs.\
13 DOF from the RS Hamiltonian is an identified open problem.
\end{remark}

\subsection{Numerical Prediction}

\begin{corollary}[RS Hubble Constant Prediction]
\label{cor:h0}
Taking $H_0^{\mathrm{Planck}} = 67.4~\mathrm{km/s/Mpc}$
as input (the CMB-derived value):
\begin{equation}
  H_0^{\mathrm{RS}} = 67.4 \times \frac{13}{12}
  = 73.0~\mathrm{km/s/Mpc}.
\end{equation}
Note: this uses the CMB measurement as an input; the prediction is
the ratio $13/12$, not the absolute value $73.0~\mathrm{km/s/Mpc}$.
\end{corollary}

\begin{remark}[Comparison with Observation]
\begin{center}
\renewcommand{\arraystretch}{1.3}
\begin{tabular}{@{}lcc@{}}
\toprule
\textbf{Quantity} & \textbf{RS prediction} &
\textbf{Observed} \\
\midrule
$H_0^{\mathrm{late}} / H_0^{\mathrm{early}}$ &
$13/12 = 1.0833$ & $73.04/67.4 = 1.084$ \\
$H_0^{\mathrm{late}}$ &
$73.0~\mathrm{km/s/Mpc}$ &
$73.0 \pm 1.0$ \\
Agreement & --- & $0.03\%$ \\
\bottomrule
\end{tabular}
\end{center}
\end{remark}

%% ============================================================
\section{Physical Interpretation}
\label{sec:interpretation}
%% ============================================================

\begin{remark}[No New Physics Required --- RS Interpretation]
\label{prop:no-new}
If the $13/12$ identification is correct, the Hubble tension is not a
discrepancy requiring new physics (early dark energy, modified
recombination, BSM particles) but a geometric artifact of the
different numbers of degrees of freedom accessed by
early-universe vs.\ late-universe measurements.
\end{remark}

\begin{remark}
The analogy: measuring the size of a room with a ruler
(12 edges) vs.\ measuring the size of a room while it is
growing (12 edges + the temporal growth direction = 13).
The late-universe ``ruler'' has one more tick mark, giving
a systematically larger rate.
\end{remark}

\begin{remark}[Consistency Checks]
Independent measurements that probe intermediate
redshifts ($z \sim 0.5$--$2$) should yield $H_0$ values
between the Planck and SH0ES values, interpolating between
the 12-DOF and 13-DOF regimes.  Recent DESI BAO results
suggest $H_0 \approx 68$--$70~\mathrm{km/s/Mpc}$,
consistent with this interpolation.
\end{remark}

%% ============================================================
\section{Dark Energy: $\Omega_\Lambda = 11/16 - \alpha/\pi$}
\label{sec:dark-energy}
%% ============================================================

\begin{proposition}[Dark Energy Density Parameter --- RS Structural Prediction]
\label{thm:omega-lambda}
The dark energy density parameter is:
\begin{equation}
  \boxed{\Omega_\Lambda =
  \frac{E_{\mathrm{passive}}}{2V} - \frac{\alpha}{\pi}
  = \frac{11}{16} - \frac{\alpha}{\pi} \approx 0.6852.}
  \label{eq:omega-lambda}
\end{equation}
(Full derivation in the companion paper on the cosmological constant~\cite{RS_Axioms}.)
\end{proposition}

\begin{proof}
\textbf{Leading term.}\quad The 11 passive edges of the
$D = 3$ cube exert a geometric stress on the expansion---they
represent the fraction of the cube's structure that resists
dynamic change.  The total phase space volume is
$2V = 2 \times 8 = 16$ (the doubled voxel count, accounting
for debit--credit pairing in the ledger).  The passive
fraction is:
\begin{equation}
  \frac{E_{\mathrm{passive}}}{2V} = \frac{11}{16}
  = 0.6875.
\end{equation}

\textbf{Radiative correction.}\quad
The fine-structure constant $\alpha$ enters through the interaction
between the passive edge structure and the electromagnetic sector of
the ledger.  Each passive edge couples electromagnetically with
amplitude $\sim\alpha$, distributed over the phase-space angle $\pi$,
giving:
\begin{equation}
  \delta\Omega = -\frac{\alpha}{\pi}
  \approx -\frac{1/137.036}{3.14159} \approx -0.0023.
\end{equation}

Combining:
\begin{equation}
  \Omega_\Lambda = 0.6875 - 0.0023 = 0.6852.
\end{equation}
\end{proof}

\begin{remark}[Status of the $-\alpha/\pi$ correction]
\label{rem:correction-status}
The $-\alpha/\pi$ correction is a \emph{structural estimate}, not a
rigorous calculation.  The physical rationale (electromagnetic coupling
to passive edges distributed over $\pi$ steradians) is suggestive but
does not constitute a derivation from the RS Hamiltonian.  The numerical
agreement $\Omega_\Lambda^{\rm RS} = 0.6852$ vs.\ Planck $0.6847 \pm
0.0073$ is within $0.07\sigma$, but the correction term $\alpha/\pi$
was not independently derived before observing this agreement.  A
rigorous derivation of the $\alpha/\pi$ correction from the RS ledger
vertex structure is an identified open problem.  The leading term
$11/16 = 0.6875$ is a zero-parameter prediction; only the correction
has ad hoc motivation.
\end{remark}

\begin{remark}[Comparison with Observation]
\begin{center}
\renewcommand{\arraystretch}{1.3}
\begin{tabular}{@{}lcc@{}}
\toprule
\textbf{Quantity} & \textbf{RS prediction} &
\textbf{Planck 2018} \\
\midrule
$\Omega_\Lambda$ & $0.6852$ &
$0.6847 \pm 0.0073$ \\
Deviation & --- & $0.07\sigma$ \\
\bottomrule
\end{tabular}
\end{center}
\end{remark}

\begin{remark}[Interpretation]
Dark energy is \emph{not} a new substance (quintessence,
cosmological constant from vacuum energy, etc.).  It is the
geometric stress of the passive edges of the $D = 3$ cube.
The 11 passive edges resist temporal advance, producing an
accelerating expansion.  The $-\alpha/\pi$ correction
captures the electromagnetic back-reaction on this geometric
stress.
\end{remark}

%% ============================================================
\section{Cosmological Budget}
\label{sec:budget}
%% ============================================================

\begin{proposition}[Cosmic Budget --- Partial RS Derivations]
\label{prop:budget}
The RS structural estimates for cosmological density parameters
(note: only $\Omega_\Lambda$ is fully derived in this paper;
$\Omega_{\mathrm{dm}}$ and $\Omega_b$ derivations are in
companion papers):
\begin{center}
\renewcommand{\arraystretch}{1.3}
\begin{tabular}{@{}lccc@{}}
\toprule
\textbf{Component} & \textbf{RS} &
\textbf{Planck 2018} & \textbf{Source} \\
\midrule
$\Omega_\Lambda$ & $0.685$ & $0.685 \pm 0.007$ &
$11/16 - \alpha/\pi$ \\
$\Omega_{\mathrm{dm}}$ & $0.265$ & $0.265 \pm 0.007$ &
12-channel interference$^\dagger$ \\
$\Omega_b$ & $0.049$ & $0.049 \pm 0.001$ &
BBN compatibility$^\dagger$ \\
$\Omega_r$ & $\sim 10^{-4}$ & $9.1 \times 10^{-5}$ &
CMB temperature \\
\midrule
$\Omega_{\mathrm{tot}}$ & $\approx 1.000$ & $1.000 \pm 0.002$
& Flatness \\
\bottomrule
\end{tabular}
\end{center}
\end{proposition}

\noindent $^\dagger$These entries are structural estimates from companion
papers; their full derivations are under development.

\begin{remark}
The consistency of the cosmic budget
$\Omega_\Lambda + \Omega_{\mathrm{dm}} + \Omega_b +
\Omega_r \approx 1$ is a non-trivial check.  The RS
predictions for $\Omega_\Lambda$ and $\Omega_{\mathrm{dm}}$
are independently derived from cube geometry and sum to
$\approx 0.950$, leaving the correct $\sim 5\%$ for
baryonic matter.
\end{remark}

%% ============================================================
\section{The $S_8$ Tension}
\label{sec:s8}
%% ============================================================

\begin{remark}[Status of $S_8$ Predictions]
The $S_8 \equiv \sigma_8 \sqrt{\Omega_m/0.3}$ tension
between CMB ($S_8 = 0.834 \pm 0.016$) and weak lensing
surveys (KiDS-1000: $S_8 = 0.766 \pm 0.020$) represents a
$2$--$3\sigma$ discrepancy.

The RS framework attributes this to \emph{recognition
strain}: the 8-tick cycle imposes a coherence scale
$\lambda_8 \approx 8\,h^{-1}$~Mpc, below which the
$J$-cost strain $Q(k) = J(k/k_8)$ suppresses structure
growth.  A complete derivation of $\sigma_8$ from the
$\varphi$-ladder structure remains under development
(see companion paper).
\end{remark}

%% ============================================================
\section{Predictions and Falsifiers}
\label{sec:predictions}
%% ============================================================

\begin{enumerate}
  \item \textbf{$H_0$ ratio}:
  $H_0^{\mathrm{SH0ES}}/H_0^{\mathrm{CMB}} = 13/12
  = 1.0833$.  Any future measurement showing a ratio
  significantly different from 13/12 would falsify this
  prediction.

  \item \textbf{Dark energy density}:
  $\Omega_\Lambda = 11/16 - \alpha/\pi \approx 0.685$.
  Precise determinations from DESI, Euclid, and the Vera
  Rubin Observatory will test this.

  \item \textbf{Constant dark energy}: The dark energy
  equation of state $w = -1$ exactly (geometric stress
  is constant).  Any detection of evolving $w(z) \neq -1$
  would falsify the passive-edge interpretation.

  \item \textbf{No early dark energy}: The Hubble tension
  is resolved without early dark energy.  Detection of
  an early dark energy component at recombination would
  falsify the 13/12 mechanism.

  \item \textbf{Intermediate $H_0$}: BAO measurements at
  $z \sim 0.5$--$2$ should yield $H_0$ values between the
  Planck and SH0ES values, interpolating between the 12
  and 13 DOF regimes.

  \item \textbf{No new relativistic species}: $N_{\mathrm{eff}}
  = 3.046$ (standard).  Detection of $\Delta N_{\mathrm{eff}}
  > 0$ at recombination would be inconsistent with the
  RS resolution.
\end{enumerate}

%% ============================================================
\section{Conclusion}
\label{sec:conclusion}
%% ============================================================

The Hubble tension is resolved by the $D = 3$ cube geometry
of the recognition ledger.  The CMB measures $H_0$ using
$E = 12$ static edge degrees of freedom; the local distance
ladder measures using $E + 1 = 13$ dynamic degrees of freedom.
The ratio $13/12 = 1.0833$ predicts
$H_0^{\mathrm{late}} = 73.0~\mathrm{km/s/Mpc}$ from
$H_0^{\mathrm{early}} = 67.4~\mathrm{km/s/Mpc}$ (agreement
to $0.03\%$).

Dark energy is the geometric stress of the 11 passive edges,
giving $\Omega_\Lambda = 11/16 - \alpha/\pi \approx 0.685$
(within $0.1\sigma$ of Planck).  Both results are
zero-parameter predictions from cube combinatorics, requiring
no new particles, fields, or modifications to gravity.

%% ============================================================
\begin{thebibliography}{99}

\bibitem{RS_Axioms}
J.~Washburn, ``The Algebra of Reality: A Recognition Science Derivation
of Physical Law,'' \emph{Axioms} \textbf{15}(2), 90 (2026).
\texttt{doi:10.3390/axioms15020090}

\bibitem{Riess2022}
A.~G.~Riess \textit{et al.}\ (SH0ES), ``A Comprehensive
Measurement of the Local Value of the Hubble Constant,''
Astrophys.\ J.\ Lett.\ \textbf{934}, L7 (2022).

\bibitem{Planck2020}
Planck Collaboration, ``Planck 2018 Results.\ VI.\
Cosmological Parameters,'' Astron.\ Astrophys.\
\textbf{641}, A6 (2020).

\bibitem{Freedman2021}
W.~L.~Freedman, ``Measurements of the Hubble Constant:
Tensions in Perspective,'' Astrophys.\ J.\ \textbf{919},
16 (2021).

\bibitem{DiValentino2021}
E.~Di Valentino \textit{et al.}, ``In the Realm of the
Hubble Tension---A Review of Solutions,'' Class.\ Quant.\
Grav.\ \textbf{38}, 153001 (2021).

\bibitem{KiDS2021}
KiDS Collaboration, ``KiDS-1000 Cosmology: Cosmic Shear
Constraints on the Amplitude of Matter Fluctuations,''
Astron.\ Astrophys.\ \textbf{645}, A104 (2021).

\bibitem{DESI2024}
DESI Collaboration, ``DESI 2024 VI: Cosmological
Constraints from the Measurements of Baryon Acoustic
Oscillations,'' arXiv:2404.03002 (2024).

\end{thebibliography}

\end{document}
